
\chapter{積分の基礎}

\section{積分の基本性質}

\subsection{積分範囲を繋ぐ}

以下の積分範囲を繋ぐような操作はよく使います。
\begin{theorem}
  $I$を区間、$f:I\to\mathbb{R}$を連続関数とする。
  さらに$a,c,b\in I$が$a<c<b$を満たすとする。
  このとき、
  \[
    \int_a^bf(x)dx=\int_a^c f(x)dx + \int_c^bf(x)dx
  \]
\end{theorem}
\begin{proof}
  まず$a,c,b\in I$が$a<c<b$を満たすとする。
  このとき$[a,b]$の分割として
  \[
    \Delta:=\{x_0,\dots,x_j,\dots,x_m\},\quad x_j=c
  \]
  となるようにとり
  \[
    \sum_{k=1}^mf(\xi_k)(\xi_k-\xi_{k-1})=\sum_{k=1}^jf(\xi_k)(\xi_k-\xi_{k-1})+\sum_{k=j+1}^mf(\xi_k)(\xi_k-\xi_{k-1})
  \]
  において、$|\Delta|\to0$の極限を考えることで
  \[
    \int_a^bf(x)dx=\int_a^c f(x)dx + \int_c^bf(x)dx
  \]
  が得られる。
\end{proof}
\begin{corollary}
  $I$を区間、$f:I\to\mathbb{R}$を連続関数とする。
  任意の$a,b\in I$に対して
  \[
    \int_a^bf(x)dx=-\int_b^af(x)dx
  \]
\end{corollary}
\begin{proof}
  \[
    \int_a^bf(x)dx+\int_b^af(x)dx=\int_a^af(x)dx=0
  \]
\end{proof}

\subsection{積分の線形性}

\begin{theorem}
  $I$を区間、$f,g:I\to\mathbb{R}$を連続関数、$a,b\in\mathbb{R}$とする。
  このとき、連続関数$af+bg:I\to\mathbb{R}$の不定積分に関して、以下が成り立つ。
  \[
    \int(af(x)+bg(x))dx=a\int f(x)dx+b\int g(x)dx
  \]
\end{theorem}
\begin{proof}
  微分積分学の基本定理と微分の線形性から、両辺の微分が一致することは明らか。
  従って
  \[
    \int(af(x)+bg(x))dx-\left\{a\int f(x)dx+b\int g(x)dx\right\}
  \]
  は定数であるが、不定積分は定数の違いは区別しないため、題意が成り立つ。
\end{proof}

\begin{example}
  多項式
  \[
    f(x)=a_nx^n+a_{n-1}x^{n-1}+\cdots+a_0
  \]
  を積分すると、
  \[
    \int f(x)dx=\frac{a_n}{n+1}x^{n+1}+\frac{a_{n-1}}{n}x^n+\cdots+a_0x+C
  \]
  となります。$C$は積分定数です。
\end{example}

\subsection{積分の大小関係}

関数の大小関係によって、積分結果の大小関係が保存されます。
\begin{theorem}
  $I$を区間、$f,g:I\to\mathbb{R}$を連続関数とする。
  $I$上常に$f(x)\leq g(x)$ならば、
  \[
    \int_a^b f(x)\leq\int_a^b g(x)dx
  \]
  また、$I$上恒等的に$f(x)=g(x)$ではない場合に限り
  \[
    \int_a^b f(x)<\int_a^b g(x)dx
  \]
\end{theorem}
\begin{proof}
  積分の線形性から、$f(x)\geq0$ならば
  \[
    \int_a^bf(x)dx\geq0
  \]
  を示せばよいが、これは定積分の定義から明らか。

  $f(x)$が常に0ではない場合、ある$c\in I$が存在して、$f(c)>0$。
  $f$は連続だから、ある$\delta>0$が存在して、$|x-c|<\delta$を満たす任意の$x$に対して$f(x)>0$。
  よって
  \[
    \int_a^bf(x)dx=\int_a^{c-\delta}f(x)dx+\int_{c-\delta}^{c+\delta}f(x)dx+\int_{c+\delta}^bf(x)dx\geq\int_{c-\delta}^{c+\delta}f(x)dx>0
  \]
\end{proof}

\subsection{微分積分学の基本定理}

前章でほったらかしていた、微分積分学の基本定理を証明しましょう。
そのために、まずは積分に関する平均値の定理を証明します。
\begin{theorem}[積分の平均値の定理]
  $I=[a,b]$を閉区間、$f:I\to\mathbb{R}$を連続関数とする。
  このとき、ある$c\in(a,b)$が存在して、
  \[
    \frac1{b-a}\int_a^bf(x)dx=f(c)
  \]
\end{theorem}
\begin{proof}
  $f$は閉区間$[a,b]$上連続なので、最大値・最小値の定理より、最大値$M=f(x_0)$と最小値$m=f(y_0)$をもつ。
  $M=m$の場合は定理は自明なので、$M\neq m$とする。
  このとき、$m\leq f(x) \leq M$なので、
  \begin{align*}
    &m(b-a)=\int_a^bmdx\leq\int_a^b f(x)dx\leq\int_a^bMdx=M(b-a)\\
    \iff&m\leq\frac1{b-a}\int_a^bf(x)\leq M
  \end{align*}
  また、$M\neq m$ゆえ、恒等的に$m=f(x)$または$M=f(x)$ではないから、
  \[
    m<\frac1{b-a}\int_a^bf(x)< M
  \]
  が成り立つ。ゆえに中間値の定理より
  \[
    \frac1{b-a}\int_a^bf(x)dx=f(c)
  \]
  となる$c\in(a,b)$が存在する。
\end{proof}

\begin{theorem}[微分積分学の基本定理]
  $I$を区間、$f:I\to\mathbb{R}$を連続関数とする。
  このとき、任意の$a\in I$に対して、$\int_a^x f(\xi)d\xi$は微分可能な$I$上の関数であり、
  \[
    \frac{d}{dx}\int_a^x f(\xi)d\xi=f(x)
  \]
  を満たす。
\end{theorem}
\begin{proof}
  \[
    F(x):=\int_a^x f(\xi)d\xi=f(x)
  \]
  とおき、
  \[
    \frac{F(x+h)-F(x)}{h}
  \]
  を評価すると、$h>0$のとき($h<0$のときも同様である)、
  \[
    F(x+h)-F(x)=\int_a^{x+h} f(\xi) d\xi+\int_a^x f(\xi)d\xi=\int_x^{x+h}f(\xi)d\xi
  \]
  である。よって、積分に関する平均値の定理より、ある$x<\xi'<x+h$が存在して、
  \[
    \frac{F(x+h)-F(x)}{h}=f(\xi')
  \]
  これに対して$h\to0$の極限を取れば題意を得る。
\end{proof}

次の定理も微分積分学の基本公式と呼ばれます。
\begin{corollary}
  $F:[a,b]\to\mathbb{R}$を$C^1$級関数とし、$f:=F'$とする。このとき
  \[
    \int_a^b f(x)dx=F(b)-F(a)
  \]
\end{corollary}
\begin{proof}
  原始関数を、適当に$c\in I$をとって
  \[
    F(x)=\int_c^xf(\xi)d\xi+C
  \]
  とおく。
  このとき、
  \[
    F(b)-F(a)=\int_c^b f(\xi)d\xi + \int_c^af(\xi)d\xi=\int_a^bf(\xi)d\xi
  \]
\end{proof}

\subsection{部分積分}

\begin{theorem}
  $I$を区間、$f,g:I\to\mathbb{R}$を可微分関数とする。
  このとき、以下が成り立つ。
  \[
    \int f(x)g'(x)dx=f(x)g(x)-\int f'(x)g(x)dx
  \]
\end{theorem}
\begin{proof}
  積の微分公式
  \[
    (fg)'(x)=f'(x)g(x)+f(x)g'(x)
  \]
  の両辺を積分すると、積分の線形性から
  \[
    f(x)g(x)+C=\int f'(x)g(x)dx+\int f(x)g'(x)dx
  \]
  これを移項して、積分定数を吸収すれば、題意を得る。
\end{proof}

証明は簡単すぎるので省きますが、定積分バージョンは次のようになります。
\begin{theorem}
  $I=[a,b]$を区間、$f,g:I\to\mathbb{R}$を可微分関数とする。
  このとき、以下が成り立つ。
  \[
    \int_a^b f(x)g'(x)dx=\{f(b)g(b)-f(a)g(a)\}-\int_a^b f'(x)g(x)dx
  \]
\end{theorem}

上の定理における$\{f(b)g(b)-f(a)g(a)\}$の部分は、
\[
  \int_a^b(fg)'(x)dx
\]
に他ならないのですが、通常通り次のように書くことも多いです。
\begin{definition}
  $I$を区間、$f:I\to\mathbb{R}$を連続関数、$a,b\in I$とする。
  $F$を$f$の原始関数とするとき、$[a,b]$上の定積分を
  \[
    \int_a^bf(x)dx=[F(x)]_a^b
  \]
  と書く。
\end{definition}

\subsection{置換積分}

\begin{theorem}
  $I,J$を区間、$u:I\to J$を可微分関数とする。
  また$f:J\to\mathbb{R}$は連続関数で、$F:J\to\mathbb{R}$をその原始関数の一つとする。
  このとき、任意の$t\in I$に対して次が成り立つ\footnote{
    $x=u(t)$で微分するか$t$で微分するかの違いを明確にするため、以下では微分に$d/dx$や$d/dt$を使用します。
  }。
  \[
    \int f(u(t))\frac{du(t)}{dt}dt=F(u(t))
  \]
\end{theorem}
\begin{proof}
  右辺の$t$による微分は、合成関数の微分から
  \[
    \frac{dF(u(t))}{dt}=f(u(t))\frac{du(t)}{dt}
  \]
  となる。
  これは明らかに左辺の$t$による微分と一致するため、与式の左辺と右辺は積分定数の違いしか存在しない。
  ゆえに題意が成り立つ。
\end{proof}
上記の主張と証明は、$\int f(u(t))\frac{du(t)}{dt}dt$および$F(u(t))$が$t\in I$の関数とみて正しい等式です。
これを$x\in J$の関数とみても正しい等式に「戻す」ためには、$u$は全単射かつ、逆関数$u^{-1}$も可微分であるという仮定を必要とします。

従って定積分版は次のようになります。
\begin{theorem}
  $I=[a,b],J=[\alpha,\beta]$を区間、$u:I\to J$を全単射可微分関数とする。
  また$f:J\to\mathbb{R}$は連続関数とする。
  このとき次が成り立つ。
  \[
    \int_a^b f(u(t))\frac{du(t)}{dt}dt=\int_\alpha^\beta f(x)dx
  \]
\end{theorem}



